\rtmtight
\begin{tblr}{width=\textwidth,colspec={|Q[c,m,2.1cm]|X[l,m]|},hlines,vlines,hspan=minimal,rowsep=1.5pt,colsep=3pt,row{1}={bg=headgray,font=\bfseries}}
\SetCell{c,m}\hfil\logo\hfil & \textbf{POLITEKNIK NEGERI MALANG}\par\textbf{JURUSAN TEKNOLOGI INFORMASI}\par\textbf{PROGRAM STUDI: D4 TEKNIK INFORMATIKA} \\
\SetCell[c=2]{l} \textbf{RENCANA TUGAS MAHASISWA (RTM) DAN RUBRIK PENILAIAN} & \\
Nama Mata Kuliah & Sistem Pendukung Keputusan \\
Kode Mata Kuliah & RTI254001 \quad \textbf{sks / Jam perminggu:} 2 SKS/4 jam \quad \textbf{Semester:} 4 \\
Dosen Pengampu & Tim Dosen Mata Kuliah \\
\end{tblr}
\assessmentblock{Tugas Studi Kasus Metode MCDM}{Case Method dan praktik mandiri}{SCPMK0706-02801, SCPMK0706-02802, SCPMK0706-02803, SCPMK0706-02804, SCPMK1005-02805}{Mahasiswa menyelesaikan studi kasus penerapan metode MCDM pada persoalan nyata menggunakan spreadsheet dan/atau pemrograman.}{Menganalisis kasus, menerapkan metode terpilih, menghitung manual dan dengan alat bantu, mendokumentasikan hasil.}{Lembar kerja Excel, laporan perhitungan, dan dokumentasi analisis.}{Ketepatan penerapan metode MCDM pada kasus & 16.8 & Metode diterapkan dengan langkah yang benar, parameter tepat, dan hasil konsisten dengan teori serta dapat direproduksi. & Metode diterapkan dengan sebagian besar langkah benar, tetapi terdapat kesalahan minor pada parameter atau perhitungan. & Metode tidak diterapkan dengan benar atau hasil tidak dapat diverifikasi. \\ \\
Keakuratan perhitungan dan analisis sensitivitas & 11.2 & Semua perhitungan akurat, analisis sensitivitas dilakukan untuk setidaknya satu variabel kunci, dan hasilnya diinterpretasikan. & Perhitungan sebagian besar benar, tetapi analisis sensitivitas terbatas atau interpretasi kurang mendalam. & Terdapat kesalahan perhitungan yang signifikan dan tidak ada analisis sensitivitas. \\ \\
Kelengkapan dokumentasi dan laporan studi kasus & 5.6 & Laporan memuat latar belakang kasus, langkah metode, hasil, dan rekomendasi yang terdokumentasi rapi. & Laporan memuat sebagian besar elemen tetapi beberapa bagian kurang lengkap atau tidak terstruktur. & Laporan tidak lengkap, tidak terstruktur, atau tidak mencerminkan proses penyelesaian kasus. \\}

\assessmentblock{Kuis 1 Konsep SPK dan MCDM Dasar}{Kuis tertulis}{SCPMK0706-02801}{Mahasiswa menjawab soal perhitungan dan konseptual terkait teori SPK, MCDM, WSM, dan WPM.}{Mengerjakan soal secara mandiri dalam waktu yang ditentukan.}{Lembar jawaban kuis.}{Pemahaman konsep SPK dan teori pengambilan keputusan & 2.5 & Konsep SPK, evolusi, dan teori keputusan dijelaskan dengan benar dan menyeluruh sesuai pertanyaan. & Sebagian besar konsep benar, tetapi ada bagian yang kurang lengkap atau kurang akurat. & Jawaban salah atau hanya menghafal definisi tanpa pemahaman mendalam. \\ \\
Ketepatan perhitungan WSM dan WPM & 2.5 & Perhitungan WSM dan WPM dilakukan dengan benar mengikuti langkah-langkah yang tepat dan menghasilkan nilai yang akurat. & Langkah-langkah perhitungan sebagian besar benar, tetapi terdapat kesalahan aritmetik minor. & Perhitungan salah atau tidak dapat menyelesaikan permasalahan yang diberikan. \\}

\assessmentblock{UTS Analisis Penerapan Metode MCDM}{UTS berbasis kasus dan perhitungan}{SCPMK0706-02801, SCPMK0706-02802, SCPMK0706-02803}{Mahasiswa menganalisis dan menyelesaikan kasus pengambilan keputusan multi-kriteria menggunakan AHP, EDAS, dan metode perankingan.}{Mengerjakan soal kasus yang diberikan secara mandiri.}{Lembar jawaban UTS.}{Ketepatan penerapan AHP untuk pembobotan kriteria & 8 & Matriks perbandingan berpasangan, consistency ratio, dan bobot akhir dihitung dengan benar dan konsisten. & Matriks perbandingan dan sebagian perhitungan benar, tetapi terdapat kesalahan pada normalisasi atau consistency check. & AHP tidak diterapkan dengan benar atau consistency ratio tidak diperiksa. \\ \\
Akurasi perankingan alternatif dengan metode terpilih & 7 & Alternatif diranking dengan metode yang tepat, langkah-langkah lengkap, dan hasil sesuai dengan bobot yang ditetapkan. & Perankingan sebagian besar benar, tetapi terdapat kesalahan pada satu atau dua langkah kritis. & Perankingan salah atau tidak dapat diselesaikan. \\ \\
Analisis dan interpretasi hasil keputusan & 5 & Hasil perankingan diinterpretasikan dengan benar, alternatif terbaik dipilih dengan alasan yang logis, dan keterbatasan metode disebutkan. & Alternatif terbaik dipilih dengan benar, tetapi interpretasi kurang mendalam atau tidak mempertimbangkan keterbatasan. & Tidak mampu menginterpretasikan hasil atau salah memilih alternatif terbaik. \\}

\assessmentblock{Kuis 2 Metode Lanjutan dan Objektif}{Kuis tertulis}{SCPMK0706-02803}{Mahasiswa menjawab soal perhitungan menggunakan metode MEREC, MARCOS, atau metode perankingan lanjutan.}{Mengerjakan soal secara mandiri.}{Lembar jawaban kuis.}{Penguasaan metode MEREC atau MARCOS dalam pembobotan/perankingan & 2.5 & Langkah-langkah metode dijalankan dengan benar dari input hingga hasil perankingan akhir. & Sebagian besar langkah benar, tetapi terdapat kesalahan pada satu tahap. & Langkah-langkah tidak diikuti dengan benar atau hasil sangat jauh dari jawaban yang benar. \\ \\
Ketepatan analisis hasil dan pemilihan alternatif & 2.5 & Alternatif terbaik dipilih berdasarkan hasil perhitungan yang benar disertai alasan yang tepat. & Alternatif dipilih dengan benar, tetapi alasan kurang kuat. & Alternatif salah dipilih atau tidak ada alasan. \\}

\assessmentblock{PBL Proyek Implementasi SPK}{Project Based Learning}{SCPMK0706-02804, SCPMK1005-02805}{Mahasiswa mengembangkan sistem DSS sederhana dengan mengkombinasikan minimal dua metode MCDM, diimplementasikan dalam spreadsheet atau program, dan didokumentasikan.}{Merancang sistem, memilih metode, mengimplementasikan, menguji, dan mendokumentasikan proyek.}{File Excel/Python, laporan proyek, dan slide presentasi.}{Ketepatan pemilihan dan integrasi metode SPK & 4.5 & Minimal dua metode dipilih secara rasional sesuai karakteristik kasus, diintegrasikan dengan benar, dan hasilnya konsisten. & Metode dipilih dan digunakan, tetapi integrasi antar metode kurang optimal atau tidak semua langkah dijalankan dengan benar. & Hanya satu metode digunakan atau kombinasi metode tidak dapat menghasilkan keputusan yang bermakna. \\ \\
Implementasi sistem DSS menggunakan spreadsheet/pemrograman & 3.9 & Sistem terimplementasi, dapat dijalankan dengan input baru, formula/kode terdokumentasi, dan hasil terverifikasi. & Sistem berjalan untuk kasus yang diberikan, tetapi fleksibilitas terbatas atau dokumentasi kurang lengkap. & Sistem tidak berjalan atau tidak dapat menghasilkan output yang benar. \\ \\
Dokumentasi dan presentasi proyek & 3 & Laporan mencakup deskripsi kasus, metode, hasil, evaluasi, dan rekomendasi dengan slide presentasi yang terstruktur. & Laporan dan slide tersedia, tetapi beberapa komponen kurang lengkap atau kurang terstruktur. & Dokumentasi sangat minimal atau presentasi tidak mencerminkan isi proyek. \\}

\assessmentblock{UAS Evaluasi dan Presentasi Proyek DSS}{UAS presentasi dan defense}{SCPMK1005-02805}{Mahasiswa mempresentasikan proyek DSS akhir, mengevaluasi kinerja metode yang digunakan, dan mempertahankan keputusan desain sistem di hadapan penguji.}{Mempresentasikan proyek, menjawab pertanyaan penguji, dan menunjukkan hasil evaluasi kinerja sistem.}{Proyek DSS final, laporan evaluasi, dan bahan presentasi.}{Kelengkapan dan keakuratan implementasi sistem DSS & 10 & Sistem DSS final mencakup seluruh komponen yang direncanakan, dapat dioperasikan dengan input baru, dan semua perhitungan terverifikasi. & Sistem sebagian besar lengkap dan berjalan, tetapi beberapa komponen atau kasus uji masih memiliki kelemahan. & Sistem tidak lengkap, tidak dapat dioperasikan, atau menghasilkan output yang salah. \\ \\
Evaluasi kinerja metode dan analisis komparasi & 8.5 & Minimal dua metrik evaluasi digunakan, hasil komparasi antar metode dianalisis secara kritis, dan kesimpulan berbasis data. & Evaluasi dilakukan dengan satu metrik, komparasi terbatas, atau analisis kurang mendalam. & Tidak ada evaluasi kinerja yang terstruktur atau analisis komparasi metode. \\ \\
Kualitas presentasi, argumentasi, dan tanya jawab & 6.5 & Presentasi terstruktur, dapat menjelaskan keputusan desain dengan alasan teknis yang kuat, dan menjawab semua pertanyaan dengan tepat. & Presentasi cukup jelas, sebagian besar pertanyaan dijawab dengan benar, tetapi argumentasi kurang dalam pada beberapa poin. & Presentasi tidak terstruktur, tidak dapat mempertahankan keputusan desain, atau tidak dapat menjawab pertanyaan dasar. \\}

\begin{tblr}{width=\textwidth,colspec={|Q[l,m,3.2cm]|X[l,m]|},hlines,vlines,hspan=minimal,row{1-3}={bg=headgray},rowsep=1.5pt,colsep=3pt}
Jadwal Pelaksanaan: & Minggu 1--17 sesuai RPS. Setiap evaluasi memiliki RTM dan rubrik multi-kriteria. \\
Lain-Lain Yang Diperlukan: & LMS, spreadsheet/repositori proyek, perangkat lunak sesuai mata kuliah, dan perangkat presentasi. \\
Pustaka: & Mengacu pustaka utama dan pendukung pada bagian RPS. \\
\end{tblr}
